\documentclass[11pt]{article}

\usepackage[margin=1in]{geometry}
\usepackage{amsmath,amssymb,amsthm,mathtools}
\usepackage{graphicx}
\usepackage{microtype}
\usepackage[hidelinks]{hyperref}
\usepackage{enumitem}

\newtheorem{theorem}{Theorem}
\newtheorem{lemma}{Lemma}
\newtheorem{remark}{Remark}
\newcommand{\A}{\mathcal A}
\newcommand{\D}{\mathcal D}
\newcommand{\M}{\mathcal M}
\newcommand{\N}{\mathcal N}
\newcommand{\T}{\mathbb T}
\newcommand{\Z}{\mathbb Z}
\newcommand{\E}{\mathbb E}
\newcommand{\one}{\mathbf 1}
\newcommand{\ran}{\operatorname{Ran}}
\newcommand{\supp}{\operatorname{supp}}
\newcommand{\clos}[1]{\overline{#1}}

\title{A Zero-Determinant Right Outer Which Is Not Left Outer}
\author{Candidate counterexample for arXiv:1304.0518\\
Run \texttt{fa\_banach\_001}, lane 0}
\date{9 August 2026}

\begin{document}
\maketitle

\begin{center}
\fbox{\parbox{0.91\textwidth}{\textbf{Status.}
Candidate counterexample, likely valid, pending specialist human review.
The construction disproves the proposed left/right equivalence already for a
bounded analytic element and, using the standard independence of the Hardy
index, for every $1\leq p\leq\infty$.}}
\end{center}

\section{Source question}

Blecher and Labuschagne ask whether the determinant-positive equivalence
between left and right outer elements remains true when the
Fuglede--Kadison determinant is zero \cite[p.~10, after Proposition~2.11]{BL13}.
Their convention, which we use throughout, is
\[
 h\text{ right outer in }H^p \iff [h\A]_p=H^p,
 \qquad
 h\text{ left outer in }H^p \iff [\A h]_p=H^p.
\]

\begin{figure}[ht]
  \centering
  \includegraphics[width=\textwidth]{figures/open_problem_crop.png}
  \caption{The open question in arXiv:1304.0518, PDF page 10.}
\end{figure}

The exact question is whether, without the hypothesis $\Delta(h)>0$, right
outerness is equivalent to left outerness.  The construction below gives a
negative answer.

\section{Proof intuition}

Let $v$ be a nonconstant operator-valued inner function.  Right multiplication
by $v$ preserves a proper closed inner subspace, so a factorization $h=dv$ is
automatically a good candidate for failure of left outerness.  If $d$ were
bounded below, however, left multiplication by $d$ could not erase the defect
space of $vH^2$, and $h$ would also fail to be right outer.

The key is to take $d$ bounded and injective but extremely small on a sequence
of independent Bernoulli coordinates.  The model space $H^2\ominus vH^2$ is a
resolvent-kernel copy of $L^2(\D)$.  Dividing one of its vectors by $d$ produces
a weighted sum of independently shifted copies of $d^{-2}$.  We choose the
tail of $d^{-2}$ so heavy that this sum diverges almost surely.  Thus no
nonzero model-space vector lies in the range of left multiplication by $d$,
and $dvH^2$ is dense.  The same $d$ is nevertheless a full-support diagonal,
so $\A d$ is dense and $[\A dv]_2=H^2v\neq H^2$.

\section{Construction and statement}

Let
\[
 (\Omega,\mu)=\bigl((0,1)^\Z,\lambda^\Z\bigr)
\]
with Bernoulli shift $\sigma$, and put $\N=L^\infty(\Omega)$.  Let
\[
 \D=\N\rtimes_\sigma\Z
\]
be the finite crossed-product von Neumann algebra, with canonical trace $\tau$
and canonical implementing unitary $u$.  Thus, up to the harmless choice of
shift orientation,
\[
 ufu^*=f\circ\sigma^{-1}\qquad(f\in\N).
\]
Let
\[
 \M=L^\infty(\T)\ \bar\otimes\ \D,
 \qquad
 \A=H^\infty(\T)\ \bar\otimes\ \D.
\]
Equivalently, $\A$ is the weak-star closure of the analytic $\D$-valued
trigonometric polynomials.  This is a finite maximal subdiagonal algebra in
$\M$, with diagonal $\D$ and conditional expectation given by the zeroth
Fourier coefficient.

Fix $0<r<1$ and set $q=r^2$.  For $\omega=(\omega_k)_{k\in\Z}$ define
\[
 d(\omega)=\exp\!\left(-\frac{1}{2\omega_0}\right),
 \qquad X=d^{-2}=\exp\!\left(\frac{1}{\omega_0}\right).
\]
Then $d\in\N$ is bounded, positive, and has support $\one$, while
\[
 \tau(\log d)=-\frac12\int_0^1\frac{dt}{t}=-\infty,
 \qquad \Delta(d)=0.
\]
Finally define
\[
 v(z)=(zu-r\one)(\one-rzu)^{-1}=b_r(zu),
 \qquad h=dv.
\]
For $|z|=1$, the operator $zu$ is unitary and the scalar Blaschke transform
$b_r$ maps the unit circle to itself.  Hence $v$ is a unitary-valued,
nonconstant element of $\A$, and $h\in\A$.

\begin{theorem}\label{thm:main}
For the finite maximal subdiagonal algebra $\A\subset\M$ above, the bounded
analytic element $h=dv$ satisfies $\Delta(h)=0$ and
\[
 [h\A]_p=H^p
 \quad\text{but}\quad
 [\A h]_p\neq H^p
 \qquad(1\leq p\leq\infty).
\]
In particular, a zero-determinant right outer need not be left outer.
\end{theorem}

\section{Proof}

We first prove the two claims in $H^2$.

\subsection{The model space of the inner factor}

Identify
\[
 H^2(\A)=\ell^2\bigl(\mathbb N_0;L^2(\D)\bigr)
\]
by Fourier coefficients.  Let $W$ be left multiplication by $zu$.  It is the
pure isometry
\[
 W(x_0,x_1,x_2,\ldots)=(0,ux_0,ux_1,\ldots)
\]
whose wandering space is the constant copy of $L^2(\D)$.  Left multiplication
by $v$ is $b_r(W)$.  The usual one-point Blaschke model-space calculation,
included here to fix the module orientation, gives
\begin{equation}\label{eq:model}
 K_v:=H^2\ominus vH^2
 =\left\{k_x:x\in L^2(\D)\right\},
 \qquad
 (k_x)_n=\sqrt{1-r^2}\,r^n u^n x.
\end{equation}
Indeed,
\[
 k_x=\sqrt{1-r^2}\,(\one-rW)^{-1}x,
 \qquad W^*k_x=rk_x,
\]
so $b_r(W)^*k_x=0$.  Conversely the Cayley transform identity for $b_r$
maps $\ker W^*$ onto $\ker b_r(W)^*$, giving all of the kernel.  Also
\[
 \|k_x\|_2^2=(1-r^2)\sum_{n\geq0}r^{2n}\|u^nx\|_2^2=\|x\|_2^2,
\]
so the parametrization is isometric.

\subsection{Right outerness}

Because $d=d^*$ is bounded, a vector $y\in H^2$ is orthogonal to $h\A$ if
and only if
\[
 \langle dva,y\rangle=0\quad(a\in\A)
 \iff dy\in H^2\ominus vH^2=K_v.
\]
Suppose $dy=k_x$.  Since $d$ has support $\one$, the Fourier coefficients of
$y$ are forced to be
\[
 y_n=\sqrt{1-r^2}\,r^n d^{-1}u^nx.
\]
Consequently, if $y\in H^2$, then
\begin{align}
 \|y\|_2^2
 &=(1-r^2)\sum_{n\geq0}q^n
   \tau\!\left(x^*u^{-n}d^{-2}u^nx\right).\label{eq:energy}
\end{align}

We show that the right-hand side is infinite whenever $x\neq0$.  Put
\[
 X_n=u^{-n}Xu^n\in L^0(\N),
 \qquad
 S_N=(1-r^2)\sum_{n=0}^Nq^nX_n.
\]
The $X_n$ are independent copies of $X$.  For $t\geq e$,
\[
 \mu\{X>t\}
 =\mu\left\{\omega_0<\frac1{\log t}\right\}
 =\frac1{\log t}.
\]
Therefore, for all sufficiently large $n$,
\[
 \sum_n\mu\{X_n>q^{-2n}\}
 =\sum_n\frac1{2n|\log q|}=\infty.
\]
The second Borel--Cantelli lemma implies that $X_n>q^{-2n}$ infinitely often
almost surely.  Along those indices, $q^nX_n>q^{-n}$, and hence
\begin{equation}\label{eq:Sinfty}
 S_N(\omega)\longrightarrow\infty
 \quad\text{for almost every }\omega.
\end{equation}

Let $E_{\N}:\D\to\N$ be the canonical trace-preserving conditional
expectation.  For $x\in L^2(\D)$ put
\[
 g=E_{\N}(xx^*)\in L^1(\N)_+.
\]
Faithfulness gives $g\neq0$ if $x\neq0$.  Using truncations for the extended
positive operators and then trace cyclicity,
\begin{align*}
 &(1-r^2)\sum_{n=0}^Nq^n
   \tau\!\left(x^*X_nx\right)
   =\tau(S_Nxx^*)
   =\int_\Omega S_N g\,d\mu.
\end{align*}
By \eqref{eq:Sinfty} and monotone convergence, the last expression tends to
$\infty$ for every $x\neq0$.  Thus \eqref{eq:energy} cannot be finite unless
$x=0$.  Then $dy=0$, and the full support of $d$ gives $y=0$.

We have proved $(h\A)^\perp=\{0\}$ in $H^2$, hence
\begin{equation}\label{eq:right}
 [h\A]_2=H^2.
\end{equation}

\subsection{Failure of left outerness}

Let $e_m=\one_{\{d\geq1/m\}}\in\D$ and
\[
 a_m=d^{-1}e_m\in\D\subset\A.
\]
Then $a_md=e_m\to\one$ in $L^2$, so $[\A d]_2=H^2$.  Right multiplication
by the unitary-valued analytic function $v$ is an isometry of $H^2$ and has
closed range.  Therefore
\begin{equation}\label{eq:left}
 [\A h]_2=[\A dv]_2=H^2v.
\end{equation}
This range is proper.  Indeed, if $\one=fv$ for some $f\in H^2$, then the
boundary identity forces $f=v^*$.  But $v$ is nonconstant (its first positive
Fourier coefficient is $(1-r^2)u$), so $v^*$ has a nonzero negative Fourier
coefficient and does not belong to $H^2$.  Hence $H^2v\neq H^2$.

Equations \eqref{eq:right} and \eqref{eq:left} prove the counterexample for
$p=2$.

\subsection{All Hardy indices and the determinant}

Blecher and Labuschagne's index-independence lemma says that, for
$1\leq p\leq\infty$, an element $f\in L^p(\M)$ satisfies
$[f\A]_p=H^p$ exactly when it is right outer in $H^1$
\cite[Lemma~4.1]{BL08}; the left-handed statement follows by passing to the
opposite algebra.  Since our $h$ is bounded, the $p=2$ conclusions therefore
hold for every $1\leq p\leq\infty$.

Finally $v$ is unitary-valued, so multiplicativity of the
Fuglede--Kadison determinant gives
\[
 \Delta(h)=\Delta(d)\Delta(v)=0.
\]
Also $\Phi(v)=-r\one$ and $\Phi(h)=-rd$, so
$\Delta(\Phi(h))=0$, as expected for the determinant-zero regime of the
source question.  This completes the proof of Theorem~\ref{thm:main}.
\qed

\section{Verification notes}

\begin{enumerate}[leftmargin=*]
\item \textbf{Hypotheses.}  The algebra
$H^\infty(\T)\bar\otimes\D$ is a standard finite maximal subdiagonal algebra,
and $h=dv$ is bounded and analytic.
\item \textbf{No hidden bounded-inverse assumption.}  The proof uses only that
$d$ is bounded and has full support.  Its inverse is an affiliated operator;
the possible failure of $d^{-1}u^nx$ to lie in $L^2$ is exactly the mechanism
measured by \eqref{eq:energy}.
\item \textbf{Left/right orientation.}  Right outerness concerns left
multiplication by $dv$ on $vH^2$ and is proved through
$H^2\ominus vH^2$.  Failure of left outerness concerns right multiplication
by $v$ and yields the distinct subspace $H^2v$.
\item \textbf{Probabilistic step.}  The shifted $X_n$ are independent because
they are functions of distinct Bernoulli coordinates.  The chosen tail makes
the Borel--Cantelli sum harmonic.
\item \textbf{Extended-positive step.}  Every use of $d^{-2}$ can be read with
spectral truncations first.  Monotone convergence then justifies both trace
cyclicity and the expectation identity.
\item \textbf{Computational dependence.}  None.  The construction and proof
are exact; no finite experiment or numerical lemma is used.
\end{enumerate}

\section{Novelty and limitations}

The bounded novelty search performed on 9 August 2026 included:
\begin{itemize}[leftmargin=*]
\item the run's registry, solution, attempt, and proof-gap indexes, plus its
local 2000--2026 arXiv source corpus;
\item exact and close-variant web/arXiv searches for ``left outer'', ``right
outer'', ``determinant zero'', ``maximal subdiagonal'', and the title/arXiv id;
\item the OpenAlex citation graph for DOI \texttt{10.4064/sm217-3-4}, containing
37 records.  The seven substantive mathematical citing works concerned
Orlicz--Hardy spaces, symmetric Hardy spaces, Haagerup Hardy spaces, invariant
subspaces, or Beurling-type theorems; none found in the search stated a
resolution or counterexample to the exact determinant-zero question.
\end{itemize}
This bounded search is evidence, not a proof of novelty.  The result should be
checked by a specialist in noncommutative Hardy spaces, with particular focus
on the model-space orientation and the extended-positive conditional-
expectation calculation.  The packet makes no claim about whether every right
outer is uniform outer; it only settles the left/right equivalence negatively.

\section{Human-review recommendation}

\textbf{Recommendation: high-priority expert verification.}  The proof is
short enough to audit locally and, if correct, gives a full negative answer to
the explicit question in the source paper.  A reviewer should first verify
\eqref{eq:model}, then audit the implication $dy=k_x\Rightarrow$
\eqref{eq:energy}, and finally confirm that the crossed-product expectation
turns the energy into $\int S_Ng$ for arbitrary $x\in L^2(\D)$.

\begin{thebibliography}{9}

\bibitem{BL08}
D. P. Blecher and L. E. Labuschagne,
\emph{Applications of the Fuglede--Kadison determinant: Szeg\H{o}'s theorem
and outers for noncommutative $H^p$},
Trans. Amer. Math. Soc. \textbf{360} (2008), 6131--6147.

\bibitem{BL13}
D. P. Blecher and L. E. Labuschagne,
\emph{Outers for noncommutative $H^p$ revisited},
Studia Math. \textbf{217} (2013), 265--287;
arXiv:1304.0518.

\end{thebibliography}

\end{document}
